\documentclass[10pt,a4paper]{amsart}
\usepackage{amsmath,amssymb,amsthm,graphicx}
\usepackage{verbatim}
\newenvironment{reference}{\comment}{\endcomment}
%\newenvironment{reference}{}{}
\newenvironment{slogan}{\comment}{\endcomment}
\newenvironment{history}{\comment}{\endcomment}

% For commutative diagrams we use Xy-pic
\usepackage[all]{xy}

% We use 2cell for 2-commutative diagrams.
\xyoption{2cell}
% Theorem environments.
%
\theoremstyle{plain}
\newtheorem{theorem}[subsection]{Theorem}
\newtheorem{proposition}[subsection]{Proposition}
\newtheorem{lemma}[subsection]{Lemma}

\theoremstyle{definition}
\newtheorem{definition}[subsection]{Definition}
\newtheorem{example}[subsection]{Example}
\newtheorem{exercise}[subsection]{Exercise}
\newtheorem{situation}[subsection]{Situation}

\theoremstyle{remark}
\newtheorem{remark}[subsection]{Remark}
\newtheorem{remarks}[subsection]{Remarks}

\numberwithin{equation}{subsection}

% Macros
%
\def\lim{\mathop{\mathrm{lim}}\nolimits}
\def\colim{\mathop{\mathrm{colim}}\nolimits}
\def\Spec{\mathop{\mathrm{Spec}}}
\def\Hom{\mathop{\mathrm{Hom}}\nolimits}
\def\Ext{\mathop{\mathrm{Ext}}\nolimits}
\def\SheafHom{\mathop{\mathcal{H}\!\mathit{om}}\nolimits}
\def\SheafExt{\mathop{\mathcal{E}\!\mathit{xt}}\nolimits}
\def\Sch{\mathit{Sch}}
\def\Mor{\mathop{\mathrm{Mor}}\nolimits}
\def\Ob{\mathop{\mathrm{Ob}}\nolimits}
\def\Sh{\mathop{\mathit{Sh}}\nolimits}
\def\NL{\mathop{N\!L}\nolimits}
\def\CH{\mathop{\mathrm{CH}}\nolimits}
\def\proetale{{pro\text{-}\acute{e}tale}}
\def\etale{{\acute{e}tale}}
\def\QCoh{\mathit{QCoh}}
\def\Ker{\mathop{\mathrm{Ker}}}
\def\Im{\mathop{\mathrm{Im}}}
\def\Coker{\mathop{\mathrm{Coker}}}
\def\Coim{\mathop{\mathrm{Coim}}}

% Boxtimes
%
\DeclareMathSymbol{\boxtimes}{\mathbin}{AMSa}{"02}

%
% Macros for moduli stacks/spaces
%
\def\QCohstack{\mathcal{QC}\!\mathit{oh}}
\def\Cohstack{\mathcal{C}\!\mathit{oh}}
\def\Spacesstack{\mathcal{S}\!\mathit{paces}}
\def\Quotfunctor{\mathrm{Quot}}
\def\Hilbfunctor{\mathrm{Hilb}}
\def\Curvesstack{\mathcal{C}\!\mathit{urves}}
\def\Polarizedstack{\mathcal{P}\!\mathit{olarized}}
\def\Complexesstack{\mathcal{C}\!\mathit{omplexes}}
% \Pic is the operator that assigns to X its picard group, usage \Pic(X)
% \Picardstack_{X/B} denotes the Picard stack of X over B
% \Picardfunctor_{X/B} denotes the Picard functor of X over B
\def\Pic{\mathop{\mathrm{Pic}}\nolimits}
\def\Picardstack{\mathcal{P}\!\mathit{ic}}
\def\Picardfunctor{\mathrm{Pic}}
\def\Deformationcategory{\mathcal{D}\!\mathit{ef}}
\usepackage{cleveref}
\crefname{originalsection}{original section}{original sections}
\Crefname{originalsection}{Original section}{Original sections}
\makeatletter
\newlabel{duality-example-iso-on-RSheafHom-noetherian}{{0GEW}{source}{}{sourcepointer-1}{}}
\newlabel{duality-example-iso-on-RSheafHom-noetherian@cref}{{[originalsection][0][]0GEW}{[1][0][]source}{}{}{}}
\makeatother
\crefname{originalsection}{original section}{original sections}
\Crefname{originalsection}{Original section}{Original sections}
\makeatletter
\newlabel{duality-example-proper-over-local}{{0AWI}{source}{}{sourcepointer-2}{}}
\newlabel{duality-example-proper-over-local@cref}{{[originalsection][0][]0AWI}{[1][0][]source}{}{}{}}
\makeatother
\crefname{originallemma}{original lemma}{original lemmas}
\Crefname{originallemma}{Original lemma}{Original lemmas}
\makeatletter
\newlabel{duality-lemma-dualizing-module-CM-scheme}{{0AWT}{source}{}{sourcepointer-3}{}}
\newlabel{duality-lemma-dualizing-module-CM-scheme@cref}{{[originallemma][0][]0AWT}{[1][0][]source}{}{}{}}
\makeatother
\crefname{originallemma}{original lemma}{original lemmas}
\Crefname{originallemma}{Original lemma}{Original lemmas}
\makeatletter
\newlabel{duality-lemma-shriek-dualizing}{{0AA3}{source}{}{sourcepointer-4}{}}
\newlabel{duality-lemma-shriek-dualizing@cref}{{[originallemma][0][]0AA3}{[1][0][]source}{}{}{}}
\makeatother
\crefname{originalsection}{original section}{original sections}
\Crefname{originalsection}{Original section}{Original sections}
\makeatletter
\newlabel{dualizing-section-dualizing-local}{{0A7M}{source}{}{sourcepointer-5}{}}
\newlabel{dualizing-section-dualizing-local@cref}{{[originalsection][0][]0A7M}{[1][0][]source}{}{}{}}
\makeatother
\crefname{originalreference}{original reference}{original references}
\Crefname{originalreference}{Original reference}{Original references}
\makeatletter
\newlabel{lemma-R1-injective}{{0AXC}{source}{}{sourcepointer-6}{}}
\newlabel{lemma-R1-injective@cref}{{[originalreference][0][]0AXC}{[1][0][]source}{}{}{}}
\makeatother
\crefname{originallemma}{original lemma}{original lemmas}
\Crefname{originallemma}{Original lemma}{Original lemmas}
\makeatletter
\newlabel{properties-lemma-normal-dimension-2-Cohen-Macaulay}{{0B3D}{source}{}{sourcepointer-7}{}}
\newlabel{properties-lemma-normal-dimension-2-Cohen-Macaulay@cref}{{[originallemma][0][]0B3D}{[1][0][]source}{}{}{}}
\makeatother
\crefname{originalreference}{original reference}{original references}
\Crefname{originalreference}{Original reference}{Original references}
\makeatletter
\newlabel{remark-dualizing-setup}{{0B4R}{source}{}{sourcepointer-8}{}}
\newlabel{remark-dualizing-setup@cref}{{[originalreference][0][]0B4R}{[1][0][]source}{}{}{}}
\makeatother
\crefname{originalreference}{original reference}{original references}
\Crefname{originalreference}{Original reference}{Original references}
\makeatletter
\newlabel{situation-vanishing}{{07UY}{source}{}{sourcepointer-9}{}}
\newlabel{situation-vanishing@cref}{{[originalreference][0][]07UY}{[1][0][]source}{}{}{}}
\makeatother
\title{Vanishing of the first direct image of the dualizing module in the rational singularity situation}
\author{The Stacks Project Authors}
\date{stacks-project:resolve.tex:local-snapshot-20261012; GNU FDL 1.2 or later}
\begin{document}
\maketitle
\noindent\textit{Editorial note.} All delivered mathematical text is the author's own, copied from the pinned Stacks Project chapter file resolve.tex; no statement or proof step was rewritten, reordered or removed.\par\medskip
\noindent\textit{Editorial note.} The counted proof is the author's complete argument: the normalization of the dualizing complex and the reduction of a failure of vanishing to a nonzero map to the residue field, the application of the dualizing functor with the two displayed identifications, and the final exclusion through the injectivity of the restriction of the first cohomology of the structure sheaf.\par\medskip
\noindent\textit{Editorial note.} The label proposition-Grauert-Riemenschneider names the counted statement, and the proof environment immediately following it in the source is delivered verbatim as its proof.\par\medskip
\noindent\textit{Editorial note.} The counted statement is carried by the author's own proposition environment in the chapter source.\par\medskip
\noindent\textit{Editorial note.} Pre-open gate: python operations/agent-authoring/tools/precard\_gate.py --source resolve.tex --label proposition-Grauert-Riemenschneider --start 1837 --end 1887 --local-id 28097 returned RESULT=OK (L28097 free, occupancy FREE, public interval clear). Pre-assembly tag self-check: canonical tag resolved and not present in reference/public-items.jsonl.\par\medskip
\noindent\textit{Editorial note.} The statement and proof are copied from the pinned chapter file; the cited inputs are kept as external pointers with their canonical Stacks tags.\par\medskip
\noindent\textit{Allowed standard prerequisites (editorial).}\begin{itemize}
\item The dualizing module of a normal surface as a shift of its dualizing complex. Properties, the lemma that a normal scheme of dimension two is Cohen-Macaulay, and Duality for Schemes, the lemma identifying the dualizing complex of a Cohen-Macaulay scheme with a shift of a module, together with the normalization of the dualizing complex of the base. Supplies the identification of the dualizing complex of the source with a shift of the dualizing module and the normalization used in the duality computation.
\item Duality identities for the pushforward of the dualizing complex. Duality for Schemes, the example identifying derived homs into the dualizing complex for the inclusion of a dualizing complex, the identity expressing the dualizing complex of the source as the exceptional inverse image of the one of the base, and the lemma that the exceptional inverse image is dualizing. Supply the two displayed identifications computing the derived hom of the pushforward of the dualizing complex with the structure sheaf pushforward.
\item Injectivity of the restriction of the first cohomology. Same chapter, the lemma asserting that the restriction map from the first cohomology of the structure sheaf to the inverse image of the open is injective. Supplies the contradiction excluding the map produced by duality.
\end{itemize}\medskip
\ifdefined\hypertarget\hypertarget{sourcepointer-1}{}\fi
\noindent\textit{Source reference.} Original section 0GEW: sources/stacks-chow/tags.txt, source lines 20491--20491. Prerequisite pointer invoked at proof line(s) 1882; it names a standard fact of a sibling chapter that the counted argument uses, and every step of the counted proof built on it is retained verbatim in the delivered body.\par\medskip
\ifdefined\hypertarget\hypertarget{sourcepointer-2}{}\fi
\noindent\textit{Source reference.} Original section 0AWI: sources/stacks-chow/tags.txt, source lines 13721--13721. Prerequisite pointer invoked at proof line(s) 1859; it names a standard fact of a sibling chapter that the counted argument uses, and every step of the counted proof built on it is retained verbatim in the delivered body.\par\medskip
\ifdefined\hypertarget\hypertarget{sourcepointer-3}{}\fi
\noindent\textit{Source reference.} Original lemma 0AWT: sources/stacks-chow/tags.txt, source lines 13731--13731. Prerequisite pointer invoked at proof line(s) 1857; it names a standard fact of a sibling chapter that the counted argument uses, and every step of the counted proof built on it is retained verbatim in the delivered body.\par\medskip
\ifdefined\hypertarget\hypertarget{sourcepointer-4}{}\fi
\noindent\textit{Source reference.} Original lemma 0AA3: sources/stacks-chow/tags.txt, source lines 12971--12971. Prerequisite pointer invoked at proof line(s) 1886; it names a standard fact of a sibling chapter that the counted argument uses, and every step of the counted proof built on it is retained verbatim in the delivered body.\par\medskip
\ifdefined\hypertarget\hypertarget{sourcepointer-5}{}\fi
\noindent\textit{Source reference.} Original section 0A7M: sources/stacks-chow/tags.txt, source lines 12885--12885. Prerequisite pointer invoked at proof line(s) 1852; it names a standard fact of a sibling chapter that the counted argument uses, and every step of the counted proof built on it is retained verbatim in the delivered body.\par\medskip
\ifdefined\hypertarget\hypertarget{sourcepointer-6}{}\fi
\noindent\textit{Source reference.} Original reference 0AXC: sources/stacks-resolve/pinned-original-source.tex, source lines 1757--1757. Prerequisite pointer invoked at proof line(s) 1867; it names a standard fact of the same chapter that the counted argument uses, and every step of the counted proof built on it is retained verbatim in the delivered body.\par\medskip
\ifdefined\hypertarget\hypertarget{sourcepointer-7}{}\fi
\noindent\textit{Source reference.} Original lemma 0B3D: sources/stacks-chow/tags.txt, source lines 13961--13961. Prerequisite pointer invoked at proof line(s) 1855; it names a standard fact of a sibling chapter that the counted argument uses, and every step of the counted proof built on it is retained verbatim in the delivered body.\par\medskip
\ifdefined\hypertarget\hypertarget{sourcepointer-8}{}\fi
\noindent\textit{Source reference.} Original reference 0B4R: sources/stacks-resolve/pinned-original-source.tex, source lines 1807--1807. Prerequisite pointer invoked at proof line(s) 1845; it names a standard fact of the same chapter that the counted argument uses, and every step of the counted proof built on it is retained verbatim in the delivered body.\par\medskip
\ifdefined\hypertarget\hypertarget{sourcepointer-9}{}\fi
\noindent\textit{Source reference.} Original reference 07UY: sources/stacks-resolve/pinned-original-source.tex, source lines 1584--1584. Prerequisite pointer invoked at proof line(s) 1839; it names a standard fact of the same chapter that the counted argument uses, and every step of the counted proof built on it is retained verbatim in the delivered body.\par\medskip
\begin{proposition}[Grauert-Riemenschneider]
\label{proposition-Grauert-Riemenschneider}
In Situation \ref{situation-vanishing} assume
\begin{enumerate}
\item $X$ is a normal scheme,
\item $A$ is Nagata and has a dualizing complex $\omega_A^\bullet$.
\end{enumerate}
Let $\omega_X$ be the dualizing module of $X$
(Remark \ref{remark-dualizing-setup}). Then $R^1f_*\omega_X = 0$.
\end{proposition}

\begin{proof}
In this proof we will use the identification $D(A) = D_\QCoh(\mathcal{O}_S)$
to identify quasi-coherent $\mathcal{O}_S$-modules with $A$-modules.
Moreover, we may assume that $\omega_A^\bullet$ is normalized, see
Dualizing Complexes, Section \ref{dualizing-section-dualizing-local}.
Since $X$ is a Noetherian normal $2$-dimensional scheme
it is Cohen-Macaulay (Properties, Lemma
\ref{properties-lemma-normal-dimension-2-Cohen-Macaulay}).
Thus $\omega_X^\bullet = \omega_X[2]$ (Duality for Schemes, Lemma
\ref{duality-lemma-dualizing-module-CM-scheme} and the
normalization in Duality for Schemes, Example
\ref{duality-example-proper-over-local}).
If the proposition is false, then we can find a nonzero map
$R^1f_*\omega_X \to \kappa$. In other words we obtain a nonzero map
$\alpha : Rf_*\omega_X^\bullet \to \kappa[1]$.
Applying $R\Hom_A(-, \omega_A^\bullet)$ we get a nonzero map
$$
\beta : \kappa[-1] \longrightarrow Rf_*\mathcal{O}_X
$$
which is impossible by Lemma \ref{lemma-R1-injective}.
To see that $R\Hom_A(-, \omega_A^\bullet)$ does what we said, first
note that
$$
R\Hom_A(\kappa[1], \omega_A^\bullet) =
R\Hom_A(\kappa, \omega_A^\bullet)[-1] =
\kappa[-1]
$$
as $\omega_A^\bullet$ is normalized and we have
$$
R\Hom_A(Rf_*\omega_X^\bullet, \omega_A^\bullet) =
Rf_*R\SheafHom_{\mathcal{O}_X}(\omega_X^\bullet, \omega_X^\bullet) =
Rf_*\mathcal{O}_X
$$
The first equality by
Duality for Schemes, Example \ref{duality-example-iso-on-RSheafHom-noetherian}
and the fact that $\omega_X^\bullet = f^!\omega_A^\bullet$
by construction, and the second equality because $\omega_X^\bullet$
is a dualizing complex for $X$ (which goes back to
Duality for Schemes, Lemma \ref{duality-lemma-shriek-dualizing}).
\end{proof}

\end{document}
